\begin{proof}[Proof of \cref{obj:lem:two-point-divergence}]
Take the constant
\[
C=\max\{8c_B,\,5\}>0,
\]
which will serve all three conclusions, and fix \(n\ge1\) large enough that \(0<h_n\le1/2\) and \(0<q_n\le1/2\).  Write \(h=h_n\), \(P_+=P_{n,+}\), and \(P_-=P_{n,-}\).

\emph{Step 1: the one-observation divergence.}  By \cref{obj:def:two-point-witness} the two laws have the same covariate law, the same propensity, and the same outcome law except on the treated active cell.  Partition the observation space into the three measurable cells
\[
s_0=\{\text{not }(X\in B_n\text{ and }A=1)\},
\qquad
s_1=\{X\in B_n,\ A=1,\ Y=1\},
\qquad
s_2=\{X\in B_n,\ A=1,\ Y=-1\}.
\]
On \(s_0\) the two laws coincide.  On the treated active cell the outcome is Bernoulli on \(\{-1,1\}\) with \(P_\sigma(Y=1\mid X\in B_n,A=1)=(1+\sigma h)/2\), so \(P_+\) is a constant multiple of \(P_-\) on each cell, with ratios
\[
c_0=1,\qquad c_1=\frac{1+h}{1-h},\qquad c_2=\frac{1-h}{1+h},
\]
and in particular \(P_+\ll P_-\) with a bounded density, so the squared deviation of the density is \(P_-\)-integrable.  Writing
\[
A_n=P_X(B_n)\,q_n
\]
for the probability of entering the informative cell, the two charged cells carry \(P_-\)-masses \(A_n(1-h)/2\) and \(A_n(1+h)/2\).  Since the density is constant on each cell of a finite partition, the chi-square divergence is the corresponding finite sum, which evaluates exactly:
\[
\chi^2(P_+\,\|\,P_-)
=
\sum_{i=0}^2\left(c_i-1\right)^2P_-(s_i)
=
\left(\frac{2h}{1-h}\right)^2\frac{A_n(1-h)}2
+\left(\frac{2h}{1+h}\right)^2\frac{A_n(1+h)}2
=
\frac{4A_nh^2}{1-h^2}.
\]
Because \(0\le h\le1/2\) gives \(1-h^2\ge3/4\ge1/2\), we have \(4/(1-h^2)\le8\), hence
\[
\chi^2(P_+\,\|\,P_-)\le 8A_nh^2 .
\]
% lean: chiSqDiv_three_cell_bound_of_restrict
Since \(P_X\) is Lebesgue measure restricted to \([0,1]\) and \(B_n=[0,c_Bh^\alpha]\cap[0,1]\), the block mass obeys \(P_X(B_n)\le c_Bh^\alpha\), so \(A_n\le c_Bh^\alpha q_n\) and
\[
\chi^2(P_+\,\|\,P_-)
 \le 8\,c_Bh^\alpha\,q_n\,h^2 .
\]
The weak-arm level satisfies \(q_n\le h^{\beta_{\alpha,\gamma}}\), by the case split in \cref{obj:def:two-point-witness}: if \(\beta_{\alpha,\gamma}=0\), then \(q_n=1/4\le1=h^0\); if \(\beta_{\alpha,\gamma}>0\), then \(q_n=h^{\beta_{\alpha,\gamma}}\).  Therefore
\[
\chi^2(P_+\,\|\,P_-)
 \le 8c_B\,h_n^{2+\alpha+\beta_{\alpha,\gamma}}
 \le C\,h_n^{2+\alpha+\beta_{\alpha,\gamma}},
\]
which is the asserted one-observation bound. % lean: twoPointWitness_one_draw_chiSq_bound

\emph{Step 2: the product experiment.}  Absolute continuity and square-integrability of the density deviation tensorize, so \(P_+^{\otimes n}\ll P_-^{\otimes n}\) and
\[
\int \left(\frac{dP_+^{\otimes n}}{dP_-^{\otimes n}}-1\right)^2 dP_-^{\otimes n}<\infty .
\]
% lean: Causalean.Stat.pi_iid_absolutelyContinuous % lean: Causalean.Stat.pi_iid_integrable_sq_dev
With \(D_{\alpha,\gamma}=2+\alpha+\beta_{\alpha,\gamma}>0\) and \(h_n=n^{-1/D_{\alpha,\gamma}}\) from \cref{obj:def:two-point-witness}, we have the exact calibration
\[
h_n^{D_{\alpha,\gamma}}=n^{-1},
\qquad\text{hence}\qquad
n\,\chi^2(P_+\,\|\,P_-)\le 8c_B .
\]
The chi-square divergence tensorizes as
\[
1+\chi^2(P_+^{\otimes n}\,\|\,P_-^{\otimes n})
 =
 \left(1+\chi^2(P_+\,\|\,P_-)\right)^n .
\]
% lean: Causalean.Stat.one_add_chiSqDiv_pi_iid_general
Using \(1+x\le e^x\) for \(x=\chi^2(P_+\|P_-)\ge0\) and the construction constant \(8c_B<\log5\),
\[
\left(1+\chi^2(P_+\,\|\,P_-)\right)^n
 \le \exp\left\{n\,\chi^2(P_+\,\|\,P_-)\right\}
 \le \exp\{8c_B\}
 < 5 .
\]
Therefore \(1+\chi^2(P_+^{\otimes n}\|P_-^{\otimes n})\le5\), that is
\[
\chi^2(P_+^{\otimes n}\,\|\,P_-^{\otimes n})\le4\le 5\le C .
\]
% lean: two_point_divergence
\end{proof}
